\documentclass[11pt,a4paper]{article}

\usepackage{amsmath}
\usepackage{amssymb}
\usepackage{amsthm}
\usepackage{booktabs}
\usepackage{longtable}
\usepackage[margin=2.7cm]{geometry}
\usepackage{xcolor}
\emergencystretch=3em
\usepackage[colorlinks=true,linkcolor=blue!45!black,citecolor=blue!45!black,urlcolor=blue!45!black]{hyperref}
\usepackage[backend=biber,style=alphabetic,maxbibnames=6,giveninits=true]{biblatex}
\addbibresource{references.bib}

% ---------------------------------------------------------------- environments
\theoremstyle{plain}
\newtheorem{theorem}{Theorem}[section]
\newtheorem{proposition}[theorem]{Proposition}
\newtheorem{lemma}[theorem]{Lemma}
\newtheorem{corollary}[theorem]{Corollary}
\theoremstyle{definition}
\newtheorem{definition}[theorem]{Definition}
\theoremstyle{remark}
\newtheorem{remark}[theorem]{Remark}
\newtheorem{hazard}[theorem]{Caveat}

% Every proof in this paper opens with an italic "Idea:" line.
\newcommand{\idea}[1]{\emph{Idea: #1}\par\smallskip}

% Confidence codes. Braced so that they survive inside optional arguments.
\newcommand{\cK}{{[\mathrm{K}]}}
\newcommand{\cP}{{[\mathrm{P}]}}
\newcommand{\cPs}{{[\mathrm{P}\!\cdot\!\mathrm{s}]}}
\newcommand{\cC}{{[\mathrm{C}]}}
\newcommand{\cH}{{[\mathrm{H}]}}

% ---------------------------------------------------------------- notation
\newcommand{\N}{\mathbb{N}}
\newcommand{\R}{\mathbb{R}}
\newcommand{\Z}{\mathbb{Z}}
\newcommand{\pn}{p_n}
\newcommand{\Fb}{\mathsf{F}}
\newcommand{\Tn}{T_n}
\newcommand{\gn}{g_n}
\newcommand{\limsupn}{\limsup_{n\to\infty}}

\title{Firoozbakht's Conjecture: A Structured Attack\\[2pt]
\large A machine-checked reduction, an unconditional finite-range architecture,\\
the closure of the Riemann-Hypothesis route, and what remains open}

\author{Noogram}
\date{25 July 2026}

\begin{document}
\maketitle

\begin{abstract}
Firoozbakht's conjecture asserts that $p_{n+1}^{1/(n+1)} < p_n^{1/n}$ for every $n \ge 1$,
i.e.\ that $n \mapsto p_n^{1/n}$ is strictly decreasing. It is open. This paper reports a
structured attack on it that neither proved nor refuted it, and states precisely what the
attack did establish. Five things. (i) The reduction of the conjecture to the prime-gap
inequality $g_n < T_n$, where $T_n := p_n(p_n^{1/n}-1)$, is \emph{machine-checked} in Lean~4
against Mathlib: the development builds with exactly one \texttt{sorryAx} dependent, namely
the conjecture itself. (ii) The maximal-gap pruning rule on which the only computationally
live verification route depends is licensed by an elementary \emph{monotone-bar principle},
which shows that the obligation the attack had ranked first --- monotonicity of the exact
bar $T$ --- is misdirected: the pruning may be run against Kourbatov's monotone surrogate
$S(x) = \log^2 x - \log x - 1.17$, whose oscillation is nil. (iii) An unconditional
finite-range theorem, resting on a single explicit estimate for $\pi(x)$ due to Dusart,
converts any first-occurrence prime-gap table into a proof of the conjecture up to that
table's reach; it reproduces the published constant $1920$ at $2^{64}$ with no tuning, and
it is accompanied by a table-free window $396\,738 \le p_n \le 777\,600$ together with a
proof that no presently known unconditional gap bound can extend that window. (iv) Four
theorems close the Riemann-Hypothesis route \emph{as a route}: the sharpest published
RH-conditional gap bound certifies the Firoozbakht inequality at exactly one index in its
range of availability, and no envelope of power type at any positive exponent can certify it
beyond finitely many indices. (v) An independent exhaustive sweep to $10^{11}$
($4\,118\,054\,812$ consecutive prime pairs) finds no violation and no near-miss, with
$\max_{n \ge 10} g_n/T_n = 0.8318$. We also state, at length, what the attack did \emph{not}
establish, including two named unrepaired defects in the underlying corpus and the fact that
the citation audit for this paper has not been run.
\end{abstract}

\tableofcontents

% ==========================================================================
\section{Introduction}
\label{sec:intro}

\subsection{The conjecture}

Let $p_1 = 2 < p_2 = 3 < p_3 = 5 < \dots$ enumerate the primes, one-indexed.

\begin{definition}[Firoozbakht's conjecture, $\Fb$]
\label{def:F}
$\Fb$ is the assertion
\begin{equation}
\label{eq:F}
p_{n+1}^{1/(n+1)} < p_n^{1/n} \qquad \text{for every } n \ge 1,
\end{equation}
equivalently: the sequence $\left(p_n^{1/n}\right)_{n \ge 1}$ is strictly decreasing.
\end{definition}

The conjecture is due to Farideh Firoozbakht (University of Isfahan, 1982) and was never
published by her \cite{firoozbakht1982unpublished}; its first appearance in print is
\cite[p.~185]{ribenboim2004little}. Its date and unpublished status are attested
independently by \cite[\S1]{kourbatov2015bounds}, \cite[\S1]{ferreira2017consequences} and
\cite{visser2019verifying}. Sun states it in the form ``the sequence $(\sqrt[n]{p_n})_{n\ge1}$
is strictly decreasing'' \cite[\S1]{sun2013sequence}; Kourbatov states the equivalent
$p_{k+1} < p_k^{1+1/k}$ \cite[\S1, eq.~(1)]{kourbatov2015bounds}. All of these are the same
statement, and we use them interchangeably.

\subsection{Why it is interesting, and why it is hard}

$\Fb$ is a statement about prime gaps in disguise. Writing $g_n := p_{n+1} - p_n$ and
$L_n := \log p_n$ (natural logarithm throughout), $\Fb$ is exactly the assertion that
$g_n$ stays below the threshold $T_n := p_n(p_n^{1/n} - 1)$ at every index
(Lemma~\ref{lem:reduction}), and $T_n = L_n^2 - L_n - 1 + o(1)$
\cite[\S5, Thm.~5]{kourbatov2015bounds}. So $\Fb$ implies a Cram\'er-scale gap bound with
the leading constant $1$:
\begin{equation}
\label{eq:necessary}
\Fb \;\Longrightarrow\; \bigl(\, g_k < L_k^2 - L_k - 1 \ \text{ for all } k > 9 \,\bigr),
\end{equation}
which is \cite[\S2, Thm.~1]{kourbatov2015bounds}; see also
\cite[Thm.~2.2]{ferreira2017consequences} and the weaker $+1$ variant in
\cite[\S1]{sun2013sequence}. In the converse direction,
\begin{equation}
\label{eq:sufficient}
\bigl(\, g_k < L_k^2 - L_k - 1.17 \ \text{ for all } k > 9,\ p_k \ge 29 \,\bigr)
\;\Longrightarrow\; \Fb ,
\end{equation}
which is \cite[\S4, Thm.~3]{kourbatov2015bounds}. Thus $\Fb$ is trapped between two
$\log^2$-scale gap bounds differing by an additive $0.17$.

That sandwich is the source of the difficulty. The best \emph{unconditional} prime-gap bound
available is $g_n = O\bigl(p_n^{0.525}\bigr)$ \cite{baker2001difference}; under the Riemann Hypothesis
the state of the art is $g_n \le \tfrac{22}{25}\sqrt{p_n}\,\log p_n$ for $p_n > 3$
\cite[\S1.2, Thm.~5]{carneiro2019fourier}. Both are \emph{powers} of $p_n$; $\Fb$ needs a
\emph{polylogarithmic} bound. The distance is square-root scale against log scale, and no
known method bridges it, even conditionally on RH. Section~\ref{sec:rh} turns this
observation into theorems.

On the other side, the standard heuristic points the other way. Cram\'er's random model, in
the corrected form Granville gives it, suggests
\[
\max_{p_n \le x}(p_{n+1}-p_n) \;\gtrapprox\; 2e^{-\gamma}\log^2 x,
\qquad 2e^{-\gamma} = 1.12292\ldots,
\]
\cite[preprint p.~12, after eq.~(20)]{granville1995cramer}, which is \emph{incompatible} with
$\Fb$ via \eqref{eq:necessary}. At least one of the two --- Firoozbakht's conjecture or the
Cram\'er--Granville heuristic --- must fail, and no present technique says which. This is
the single most important framing sentence of the subject, and we return to it in
\S\ref{sec:tension}.

\subsection{Literature}
\label{sec:lit}

The relevant literature falls into five groups.

\paragraph{Statement, strengthenings and consequences.}
Ribenboim records the conjecture \cite[p.~185]{ribenboim2004little}; Sun uses it as the
motivating contrast for a provable analogue about sums of primes
\cite[Thm.~2.1]{sun2013sequence}. Ferreira and Mariano derive a systematic list of
consequences \cite{ferreira2017consequences}: the necessary gap bound
(\cite[Thm.~2.2]{ferreira2017consequences}), $g_n < \sqrt{n}$ for $n \ge 3645$
(\cite[Lem.~3.2]{ferreira2017consequences}), a Sierpi\'nski-type consequence about primes in
rows of the $n\times n$ array of $1,\dots,n^2$ (\cite[Cons.~3.3]{ferreira2017consequences}),
and --- most importantly for us --- the only unconditional positive statement in the
literature about $\Fb$ itself, Theorem~\ref{thm:io} below. Visser
\cite{visser2019verifying} states and verifies the two known strengthenings, Nicholson's
(traced to an OEIS comment \cite{oeisA182514}, otherwise unpublished) and Farhadian's
\cite{farhadian2017new}, and records the implication chain
Farhadian $\Rightarrow$ Nicholson $\Rightarrow$ Firoozbakht
\cite[Conj.~2]{visser2019verifying}, proved in
\cite[Thm.~4.5]{ferreira2017consequences} together with Forgues' weaker variant.

\paragraph{Criteria.}
The necessary and sufficient conditions \eqref{eq:necessary} and \eqref{eq:sufficient} are
Kourbatov's \cite[\S2 Thm.~1, \S4 Thm.~3]{kourbatov2015bounds}, together with a family of
sufficient conditions with $b \to 1$ \cite[\S4, Thm.~4]{kourbatov2015bounds} that presuppose
the finite verification (they are stated for $p_k > 4\cdot10^{18}$). Both rest on Axler's
effective bounds for $\pi(x)$ \cite{axler2014newbounds}; see \S\ref{sec:limits} for why that
dependency is load-bearing and unresolved here.

\paragraph{Verification.}
Kourbatov verified $\Fb$ for all $p_k < 4\cdot 10^{18}$
\cite[\S4, Thm.]{kourbatov2015verification}, using the first-occurrence prime-gap table of
Oliveira e Silva, Herzog and Pardi \cite{oliveira2014goldbach}; the 2023 endnotes to
\cite{kourbatov2015verification} extend the range to $2^{64} \approx 1.8447\cdot10^{19}$ and
record that ``prime gaps of size $g < 1920$ cannot violate'' the conjecture. Visser
independently reports verification up to the 81st maximal prime gap, ``certainly for all
primes $p < 2^{64}$'' \cite[Abstract, \S1]{visser2019verifying}. We reconstruct the
architecture of this verification from first principles in \S\ref{sec:range}.

\paragraph{The heuristic side.}
Cram\'er proves $\limsup (P_{n+1}-P_n)/(\log P_n)^2 = 1$ with probability~1 for his urn
model and only \emph{suggests} the analogue for the primes
\cite[pp.~24, 26--27]{cramer1936order}; this distinction is not a pedantry and we preserve it
everywhere. Granville corrects the model and obtains the opposite prediction
\cite[preprint pp.~10, 12]{granville1995cramer}. The empirical record of the
Cram\'er--Shanks--Granville ratio $(q-p)/\log^2 p$ is A111943 \cite{oeisA111943}; Shanks
conjectured the ratio never reaches~1 \cite{shanks1964maximal}, Granville the opposite. The
record stands at $0.9206$ at $p = 1\,693\,182\,318\,746\,371$ (gap $1132$, B.~Nyman, 1999)
\cite{oeisA111943}.

\paragraph{Large gaps.}
The best lower bound on the largest gap below $X$ is
\[
G(X) \;\ge\; f(X)\,\frac{\log X \,\log\log X \,\log\log\log\log X}{(\log\log\log X)^2},
\qquad f(X)\to\infty ,
\]
\cite{ford2016large} --- an iterated-log factor above $\log X$, but still a
full power of $\log$ below the $\log^2 X$ a counterexample to $\Fb$ would need.

\subsection{What this paper claims, and the vocabulary of its claims}
\label{sec:vocab}

This paper is an honest report on an attack that did not settle the conjecture. To keep the
strength of each statement legible we fix the following vocabulary and use it without
exception.

\begin{center}
\begin{tabular}{@{}l p{0.80\textwidth}@{}}
\toprule
Code & Meaning \\
\midrule
$\cK$ & \textbf{Machine-checked} by the Lean~4 kernel against Mathlib (\S\ref{sec:lean}). \\
$\cP$ & Paper proof, derived in this work, independently re-derived by an adversarial review leg. \\
$\cPs$ & As $\cP$, but resting on a cited source not opened in this work. \\
$\cC$ & Finite computation, exhaustive, independently reproduced by two code paths. \\
$\cH$ & Heuristic. Not evidence for or against anything, and never used in a proof. \\
\bottomrule
\end{tabular}
\end{center}

The word \emph{proved}, unqualified, is reserved in this paper for statements carrying
$\cK$. Everything at $\cP$ or $\cPs$ is a paper
proof and is labelled as such at the point of use. In particular:

\begin{center}
\fbox{\parbox{0.93\textwidth}{
\textbf{$\Fb$ is neither proved nor refuted in this paper, and nothing below should be read
as evidence that either is close.} The defensible summary sentence is:
\emph{Firoozbakht's conjecture is numerically robust over the verified range and
simultaneously incompatible with the standard Cram\'er--Granville heuristic; at least one of
the two must fail, and no current technique can say which.}
}}
\end{center}

\subsection{Contributions}

\begin{enumerate}
\item \textbf{$\cK$} A Lean~4/Mathlib formalisation in which the reduction
  $\Fb \leftrightarrow (\forall n \ge 1,\, g_n < T_n)$ is machine-checked, with an
  exhaustive list-free axiom audit over 60 declarations finding exactly one
  \texttt{sorryAx} dependent --- the conjecture itself (\S\ref{sec:lean}).
\item \textbf{$\cP$} The monotone-bar principle (Lemma~\ref{lem:M}) and its
  truncated form, which discharge the practical content of the attack's ranked-first open
  obligation without proving that obligation, and show it to be misdirected
  (\S\ref{sec:monotone}).
\item \textbf{$\cPs$, $\cC$} An unconditional
  finite-range theorem (Theorem~\ref{thm:range}) resting on a single explicit $\pi(x)$
  estimate, reproducing the published constant $1920$ at $2^{64}$; plus a table-free window
  and a proof that it closes permanently (\S\ref{sec:range}).
\item \textbf{$\cP$} Four theorems closing the RH route as a route
  (\S\ref{sec:rh}), including a Lambert-$W$ determination of the critical constant $2/e$ and
  an explicit counter-model showing that Cram\'er's $\limsup$ hypothesis does not entail
  $\Fb$ over integer sequences.
\item \textbf{$\cC$} An independent exhaustive verification to $10^{11}$ with a
  verification discipline that raises rather than returns when a margin lands inside its
  error budget (\S\ref{sec:computation}).
\item An explicit account of what is \emph{not} established, including two unrepaired
  defects in the underlying corpus and the absent citation audit (\S\ref{sec:limits}).
\end{enumerate}

% ==========================================================================
\section{Setup}
\label{sec:setup}

\subsection{Notation}

Throughout, $p_n$ is the $n$-th prime with $p_1 = 2$ (one-indexed), $\pi(x)$ is the
prime-counting function, so that
\begin{equation}
\label{eq:index}
\pi(p_n) = n ,
\end{equation}
and
\[
L_n := \log p_n, \qquad g_n := p_{n+1} - p_n, \qquad
T_n := p_n\bigl(p_n^{1/n} - 1\bigr) = p_n\bigl(e^{L_n/n} - 1\bigr).
\]
All logarithms are natural. We write $\rho_n := g_n/T_n$ for the \emph{normalised objective};
$\Fb$ at $n$ is exactly $\rho_n < 1$.

Identity \eqref{eq:index} looks like bookkeeping and is not: it is the hinge on which every
analytic estimate in this paper turns, because $T_n$ depends on the pair $(p_n, n)$ and the
only way to control $n$ analytically is through $\pi$.

\subsection{The reduction}

\begin{lemma}[Gap form of $\Fb$; $\cK$]
\label{lem:reduction}
For every $n \ge 1$,
\[
p_{n+1}^{1/(n+1)} < p_n^{1/n}
\quad\Longleftrightarrow\quad
p_{n+1}^{\,n} < p_n^{\,n+1}
\quad\Longleftrightarrow\quad
g_n < T_n .
\]
\end{lemma}

\begin{proof}
\idea{Raise to the power $n(n+1)$, take logarithms, and subtract $p_n$ from both sides; every
step is an equivalence because the maps involved are strictly increasing.}
All quantities are real and exceed $1$. The map $t \mapsto t^{n(n+1)}$ is strictly increasing
on $(0,\infty)$, so raising $p_{n+1}^{1/(n+1)} < p_n^{1/n}$ to that power gives
$p_{n+1}^{\,n} < p_n^{\,n+1}$, and the step is reversible. Applying the strictly increasing
$\log$ gives $n \log p_{n+1} < (n+1)\log p_n$, i.e.\ $\log p_{n+1} < L_n(1 + 1/n)$, i.e.\
$p_{n+1} < p_n e^{L_n/n}$, i.e.\ $p_n + g_n < p_n + p_n(e^{L_n/n}-1)$. Subtracting $p_n$
gives $g_n < T_n$.
\end{proof}

Lemma~\ref{lem:reduction} is the one result in this paper carrying $\cK$: it is
formalised and machine-checked in Lean~4 (\S\ref{sec:lean}), together with the intermediate
equivalence, and both equivalences are pinned to numerals ($g_1 = 1$, $T_1 = 2^2 - 2$) so
that a mis-stated $g$ or $T$ fails to compile rather than passing quietly.

\begin{remark}[Strict versus non-strict]
\label{rem:strict}
Visser states the conjecture with $\le$ \cite[Conj.~3, eq.~(2.4)]{visser2019verifying}. The
two forms coincide if and only if $T_n$ is never an integer. That is true for $n \ge 2$:
$p_n^{1/n}$ rational would force it to be an integer and $p_n$ to be a perfect $n$-th power,
which unique factorisation forbids for a prime. We nonetheless use the strict form
throughout and never silently convert between them.
\end{remark}

\subsection{The two known unconditional facts about $\Fb$}

\begin{theorem}[Verified range; $\cC$, cited]
\label{thm:verified}
$\Fb$ holds for every $n$ with
$p_n < 2^{64} = 18\,446\,744\,073\,709\,551\,616 \approx 1.8447\cdot10^{19}$.
\end{theorem}

This is \cite[endnotes, 5 Jan 2023]{kourbatov2015verification}, corroborated by
\cite[Abstract, \S1 eq.~(1.4)]{visser2019verifying}. Its computational input is the
first-occurrence gap table of \cite{oliveira2014goldbach}. \emph{This paper does not
re-establish Theorem~\ref{thm:verified}}; \S\ref{sec:range} reconstructs the implication
from a table to the range, and \S\ref{sec:limits} records that the table itself was never
opened in this work.

\begin{theorem}[Infinitely often; cited]
\label{thm:io}
There are infinitely many $n \in \N$ with $p_n^{1/n} > p_{n+1}^{1/(n+1)}$.
\end{theorem}

This is \cite[Thm.~5.2]{ferreira2017consequences}, proved unconditionally there via Zhang's
bounded-gaps theorem (cited as \cite[Thm.~5.1]{ferreira2017consequences}). It is the only
unconditional theorem in the literature that says something positive about $\Fb$ itself.

\begin{hazard}
Theorem~\ref{thm:io} says \emph{infinitely often}, not \emph{eventually}. It is fully
compatible with infinitely many failures. It must never drift into ``$\Fb$ holds for large
$n$''.
\end{hazard}

% ==========================================================================
\section{The monotone-bar principle and the maximal-gap reduction}
\label{sec:monotone}

\subsection{The obligation, and why it is the wrong one}

Every practical verification of $\Fb$ over a range checks \emph{record} gaps rather than all
indices. The standard justification for that pruning is the following statement, which the
present attack initially ranked as its most tractable open obligation:

\begin{quote}
\textbf{(P6$'$)} For $m < n$ straddling a record gap, $T_m \le T_n$.
\end{quote}

(P6$'$) is a statement about the oscillation of the exact bar $T$, and $T$ is \emph{not}
monotone: over the primes below $3\cdot10^6$ it decreases at about $55.9\%$ of steps
(\S\ref{sec:computation}). The main finding of this section is that no consumer of (P6$'$)
needs it.

\begin{lemma}[Monotone-bar principle; $\cP$]
\label{lem:M}
Let $B$ be a real-valued function of the prime $p$, nondecreasing in $p$. Call $k$ a
\emph{$B$-breach} if $g_k \ge B(p_k)$. If a $B$-breach exists, then the least one, $k_0$,
satisfies $g_m < g_{k_0}$ for every $m < k_0$; that is, $k_0$ is a strict maximal-gap index.
\end{lemma}

\begin{proof}
\idea{Minimality gives an upper bound at every earlier index; monotonicity moves that bound
forward to $k_0$; the two chain.}
Let $m < k_0$. Then $p_m < p_{k_0}$, so $B(p_m) \le B(p_{k_0})$ by monotonicity. By
minimality of $k_0$, the index $m$ is not a $B$-breach, so $g_m < B(p_m)$. Chaining,
$g_m < B(p_m) \le B(p_{k_0}) \le g_{k_0}$.
\end{proof}

Only monotonicity of $B$ is used: no arithmetic, no analytic input, no property of the primes
beyond $p_m < p_{k_0}$. The conclusion is strict, which is exactly the difference between
``the first failure carries a record gap'' and ``the first failure carries a gap at least as
large as all earlier ones''. The lemma says nothing about the \emph{second} breach, and it
should not: breaches after the first need not be records.

Explicit bars have validity ranges, and the first few primes sit outside them, so we need a
truncated form.

\begin{lemma}[Truncated monotone-bar principle; $\cP$]
\label{lem:Mprime}
Let $B$ be nondecreasing and let $1 \le N_1 \le N_2$. Suppose
\begin{enumerate}
\item[\textup{(i)}] $g_m < B(p_m)$ for every maximal-gap index $m$ with $N_1 \le m \le N_2$;
\item[\textup{(ii)}] $\max\{g_m : m < N_1\} < B(p_{N_1})$.
\end{enumerate}
Then $g_k < B(p_k)$ for every $k$ with $N_1 \le k \le N_2$.
\end{lemma}

\begin{proof}
\idea{Run Lemma~\ref{lem:M} inside the window, and let hypothesis \textup{(ii)} absorb
everything below it in one step.}
Suppose not, and let $k_0$ be least in $[N_1,N_2]$ with $g_{k_0} \ge B(p_{k_0})$. Take any
$m < k_0$. If $m \ge N_1$, minimality gives $g_m < B(p_m) \le B(p_{k_0}) \le g_{k_0}$. If
$m < N_1$, then by (ii) and monotonicity
$g_m \le \max\{g_j : j < N_1\} < B(p_{N_1}) \le B(p_{k_0}) \le g_{k_0}$. Either way
$g_m < g_{k_0}$, so $k_0$ is a maximal-gap index in $[N_1,N_2]$ breaching $B$, contradicting
(i).
\end{proof}

\subsection{Instantiating at Kourbatov's monotone surrogate}

Put
\[
S(x) := \log^2 x - \log x - 1.17 .
\]

\begin{lemma}[$\cP$]
\label{lem:Smono}
$S$ is strictly increasing on $x > e^{1/2}$, hence at every prime.
\end{lemma}

\begin{proof}
\idea{One derivative.}
$dS/dx = (2\log x - 1)/x > 0$ for $\log x > 1/2$, and $e^{1/2} = 1.6487\ldots < 2$.
\end{proof}

\begin{lemma}[$S$ lies below the true bar; $\cPs$]
\label{lem:S2}
For every $n$ with $p_n \ge 2\,634\,800\,823$ we have $T_n > S(p_n)$.
\end{lemma}

\begin{proof}
\idea{Replace the index $n$ by an upper bound on $\pi(p_n)$, then use $e^u - 1 > u$.}
By \eqref{eq:index}, $n = \pi(p_n)$. Axler's Corollary~3.5, in the corrigendum-corrected
form, states $\log x - 1 - 1.17/\log x < x/\pi(x)$ for $x \ge 2\,634\,800\,823$
\cite{axler2014newbounds}. With $x = p_n$ and $\ell = L_n$ this gives
$\pi(x) < x/(\ell - 1 - 1.17/\ell)$, hence $u := \ell/\pi(x) > (\ell^2 - \ell - 1.17)/x$, and
therefore $T_n = x(e^u - 1) > x u > \ell^2 - \ell - 1.17 = S(p_n)$.
\end{proof}

\begin{hazard}[The dependency, named]
\label{haz:axler}
Lemma~\ref{lem:S2} rests on \cite{axler2014newbounds}, which \emph{was not opened in this
work}; it is quoted through Kourbatov's proofs. Axler's own corrigendum moved the validity
range of Corollary~3.5 from $x \ge 5.43$ to $x \ge 2\,634\,800\,823$ --- nine orders of
magnitude --- so the range is load-bearing, not decorative: in-run data below $3\cdot10^6$
exhibits $14\,426$ indices with $T_n \le S(p_n)$. Consequently
Lemma~\ref{lem:S2} and Theorem~\ref{thm:B} carry
$\cPs$, not $\cP$. See \S\ref{sec:limits} and
Remark~\ref{rem:dusartonly} for the Dusart-only fallback.
\end{hazard}

\begin{theorem}[Maximal-gap reduction; $\cPs$]
\label{thm:B}
Let $N_2 \ge 10$ and suppose $g_m < S(p_m)$ for every maximal-gap index $m \in [10, N_2]$.
Then $g_k < S(p_k)$ for every $k \in [10, N_2]$, and consequently $\Fb$ holds at every
$k \in [10, N_2]$ with $p_k \ge 2\,634\,800\,823$.
\end{theorem}

\begin{proof}
\idea{Lemma~\ref{lem:Mprime} with the monotone surrogate $S$ in place of the oscillating
$T$; the small-index hypothesis is one finite comparison.}
Apply Lemma~\ref{lem:Mprime} with $B = S$ (nondecreasing by Lemma~\ref{lem:Smono}),
$N_1 = 10$ and $N_2$ as given. Hypothesis~(i) is the assumption. Hypothesis~(ii) is the
finite fact $\max\{g_j : j \le 9\} = 6 < 6.80139\ldots = S(29) = S(p_{10})$, verified in
exact arithmetic (\S\ref{sec:computation}, item V2). Hence $g_k < S(p_k)$
throughout $[10,N_2]$. Where $p_k \ge 2\,634\,800\,823$, Lemma~\ref{lem:S2} gives
$S(p_k) < T_k$, so $g_k < T_k$, which is $\Fb$ at $k$ by Lemma~\ref{lem:reduction}.
\end{proof}

\begin{remark}[The obligation was misdirected]
\label{rem:misdirected}
Theorem~\ref{thm:B} discharges the entire practical content of (P6$'$) without proving
(P6$'$). The oscillation of $T$ is irrelevant to the pruning, because the pruning never runs
against $T$: it runs against $S$, which does not oscillate. The pieces were all present in
the attack's own concept-card set --- the bar $S$, the (unstated but immediate) monotonicity
of $S$, and the open obligation --- and were never joined. The join is
Lemma~\ref{lem:Mprime}.
\end{remark}

\begin{remark}[Range discipline]
Theorem~\ref{thm:B} delivers nothing below $p_k = 2\,634\,800\,823$, where
Lemma~\ref{lem:S2} is unavailable and in fact false. That segment is covered by direct
computation, i.e.\ by Theorem~\ref{thm:verified}, and the two must be joined explicitly,
never assumed to overlap.
\end{remark}

\begin{remark}[The Dusart-only fallback]
\label{rem:dusartonly}
The Axler dependency of Lemma~\ref{lem:S2} can be removed entirely at the cost of a looser
bar: Lemma~\ref{lem:floor} below proves $T_n > B(p_n)$ with $B(x) = \log^2 x - 1.1\log x$
for $p_n \ge 60\,184$, from Dusart alone, and $B$ is monotone by Lemma~\ref{lem:Bmono}.
Substituting $B$ for $S$ throughout Theorem~\ref{thm:B} gives a fully $\cP$ version,
with a bar looser by roughly $0.1 L$ --- a relative deficit measured at $0.193\%$ at
$p = 2\cdot10^7$ and $0.196\%$ at $p = 1.69\cdot10^{15}$. We state the Axler version because
it is the one the literature uses; we flag the fallback because it is the one that needs no
unopened source.
\end{remark}

\subsection{Where the exact bar still resists}

Lemma~\ref{lem:M} does not settle the statement ``if $\Fb$ fails, its first failure carries
a record gap'' against the \emph{exact} bar $T$, because $T$ is not monotone. What can be
proved unconditionally is a windowed version.

\begin{theorem}[Near-record maximality, Dusart only; $\cP$]
\label{thm:C}
Suppose $\Fb$ is false and let $n_0$ be its least failure. By Theorem~\ref{thm:verified},
$p_{n_0} > 2^{64}$, hence $L_{n_0} > 44.36$ and $T_{n_0} > 1919$. Let $m < n_0$ with
\[
p_m \;\le\; p_{n_0}\, e^{-0.0623} \;=\; 0.93960 \cdots p_{n_0} .
\]
Then $g_m < g_{n_0}$.
\end{theorem}

\begin{proof}
\idea{Bound $T_m$ above and $T_{n_0}$ below with Dusart's two-sided $\pi(x)$ estimates and
ask how far apart the two logarithms must be for the upper bound at $m$ to fall below the
lower bound at $n_0$; the answer is a constant.}
If $p_m < 60\,184$ then $g_m \le 72$ (the largest gap below $60\,184$; \S\ref{sec:computation},
item V1) while $g_{n_0} \ge T_{n_0} > 1919$, and we are done. Otherwise
$p_m \ge 60\,184 > 5393$, so both halves of \cite[Thm.~6.9, eq.~(6.6)]{dusart2010estimates}
apply. Put $\ell = L_m \ge \log 60\,184 = 11.005$, $\lambda = L_{n_0} = \ell + d$ and
$\varepsilon := (\ell^2-\ell)^2/p_m$. It suffices that
$(\ell^2-\ell)\bigl(1 + (\ell^2-\ell)/p_m\bigr) \le \lambda^2 - 1.1\lambda$, i.e.\ that
\[
2\ell d + d^2 - 1.1 d - 0.1\ell \;\ge\; \varepsilon .
\]
The left side is increasing in $d \ge 0$, so it suffices that
$d \ge d^*(\ell) := (0.1\ell + \varepsilon)/(2\ell - 1.1)$. Since
$p_m \ge \max(e^{\ell}, 60\,184) = e^{\ell}$ for $\ell \ge \ell_0 := \log 60\,184 = 11.00516$,
we have $\varepsilon(\ell) \le (\ell^2-\ell)^2 e^{-\ell}$ there; that majorant has
logarithmic derivative $2(2\ell-1)/(\ell^2-\ell) - 1$, equal to $-0.618 < 0$ at $\ell_0$ and
decreasing thereafter, so $\varepsilon(\ell) \le \varepsilon(\ell_0) = 0.20145$. Hence
$d^*(\ell) \le (0.1\ell + 0.20145)/(2\ell - 1.1)$, whose numerator derivative test
$0.1(2\ell-1.1) - 2(0.1\ell + 0.20145) = -0.5129 < 0$ shows it is decreasing, so it is
maximised at $\ell_0$, where it equals $0.062264$. Thus $d \ge 0.0623$ suffices for every
admissible $\ell$.
\end{proof}

\begin{remark}[The Axler-sharpened version is withheld]
\label{rem:C-b}
A companion statement sharpens the exponent $0.0623$ by roughly a factor of ten under
Axler's corollaries. \textbf{We deliberately do not state its constant.} An adversarial
review of the underlying corpus found that the sharpened version's derivation misapplies its
upper bound on $T_m$ by a factor of order $\ell^2$, and that the in-run numerical check
reproduced the error rather than catching it --- a check written from a derivation cannot
catch an error in the derivation. The review also established that the sharpened statement's
\emph{conclusion} survives under the corrected form of the bound, but the repair has not been
applied to the source document. This paper therefore quotes only
Theorem~\ref{thm:C}, whose constant was independently confirmed conservative and which uses
no unopened source. See \S\ref{sec:limits}.
\end{remark}

\subsection{The obstruction inside the residual window}
\label{sec:obstruction}

Theorem~\ref{thm:C} leaves a sliver: the primes $p_m$ within a factor $e^{-0.0623}$ of
$p_{n_0}$. The obstruction there is exact and worth naming, because it explains why no
sharpening of constants closes the window.

To first order in $y/p_m$, with $y = p_n - p_m$ and $k = \pi(p_n) - \pi(p_m)$,
\[
T_m \le T_n
\quad\Longleftrightarrow\quad
k \;\le\; y\Bigl(1 + \frac{1}{L_m}\Bigr)\frac{\pi(p_m)}{p_m}\bigl(1 + O(y/p_m)\bigr) ,
\]
i.e., since $\pi(x)/x \approx 1/(\log x - 1)$: \emph{$T_m \le T_n$ holds exactly when the
interval $(p_m, p_n]$ contains no more than about $y/(L-2)$ primes.} The surplus
$1/(L-1) - 1/L \approx 1/L^2$ between the running-average density and the local density is
the entire drift that makes $T$ coarsely increase.

\begin{proposition}[Obstruction; $\cP$]
\label{prop:obstruction}
Closing the residual window of Theorem~\ref{thm:C} requires an unconditional upper bound on
the number of primes in an interval of length $y$ of order $x/\log x$ that is within a factor
$1 + 2/L$ of the truth. The unconditionally available bound of Brun--Titchmarsh type gives
only a factor of about $2$.
\end{proposition}

The required relative slack over the prime-number-theorem count is $4.7\%$ at $2^{64}$ and
shrinks to $1\%$ at $L = 200$, while the Brun--Titchmarsh-to-needed ratio stays near $2.1$
across that range. So the needed estimate is not a constant improvement but an
asymptotically exact short-interval count. \emph{This is, in our judgement, the most
tractable genuinely open analytic node the attack isolated --- and it is still an open
problem in analytic number theory.}

\begin{hazard}[Unsourced, and used only negatively]
\label{haz:bt}
Neither the Brun--Titchmarsh inequality nor the second Hardy--Littlewood conjecture is
covered by a source-ledger row in this work. Both are standard and undisputed; they are used
here only to \emph{name} an obstruction, never to support a positive claim.
Proposition~\ref{prop:obstruction} is stated at $\cP$ for its criterion
($1+2/L$, derived above) and its comparison to Brun--Titchmarsh is flagged as unsourced.
\end{hazard}

% ==========================================================================
\section{The unconditional finite-range architecture}
\label{sec:range}

The published frontier $2^{64}$ is not a per-pair sweep: it is a theorem about a table. This
section reconstructs that theorem from first principles, with explicit constants, resting on
\emph{one} cited inequality --- and in particular avoiding the Axler chain of
Caveat~\ref{haz:axler}.

\subsection{The analytic input}

\begin{lemma}[Dusart's upper bound for $\pi$; cited, tier L0]
\label{lem:dusart}
$\pi(x) \le x/(\log x - 1.1)$ for all $x \ge 60\,184$.
\end{lemma}

This is \cite[Thm.~6.9, eq.~(6.6)]{dusart2010estimates}, read at the locator. It is the only
external mathematical input to Theorem~\ref{thm:range}.

\begin{lemma}[Explicit floor under the bar; $\cP$]
\label{lem:floor}
Set $B(x) := \log^2 x - 1.1\log x$. Then $T_n > B(p_n)$ for every $n$ with
$p_n \ge 60\,184$.
\end{lemma}

\begin{proof}
\idea{The only opaque ingredient of $T_n$ is the index $n = \pi(p_n)$; an \emph{upper} bound
on $\pi$ turns the bar into a lower bound in $\log p_n$ alone.}
Since $e^t - 1 > t$ for $t > 0$ and $L_n/n > 0$, we have $T_n > p_n L_n/n$. By
\eqref{eq:index}, $n = \pi(p_n)$, and Lemma~\ref{lem:dusart} at $x = p_n \ge 60\,184$ gives
$n \le p_n/(L_n - 1.1)$, with $L_n - 1.1 > 0$ in this range. Hence $p_n/n \ge L_n - 1.1$ and
\[
T_n \;>\; L_n \cdot \frac{p_n}{n} \;\ge\; L_n(L_n - 1.1) \;=\; B(p_n) . \qedhere
\]
\end{proof}

The direction is the content. $T_n$ \emph{decreases} in $n$ at fixed $p_n$, so lower-bounding
$T_n$ requires an upper bound on the rank, which is precisely what Dusart's upper bound on
$\pi$ supplies. The validity range is load-bearing and not decoration: $L(L-1.1)$ crosses
$L^2 - L - 1$ --- the true asymptotic size of $T_n$ --- exactly at $L = 10$, i.e.\
$x = e^{10} = 22\,026$, below which Lemma~\ref{lem:floor} is false. Dusart's threshold
$60\,184$ clears that crossing, but only by about $0.10$ in $L$.

\begin{lemma}[$\cP$]
\label{lem:Bmono}
$B$ is strictly increasing on $[e^{0.55},\infty) \supseteq [2,\infty)$.
\end{lemma}

\begin{proof}
\idea{One derivative, in the variable $L$.}
With $L = \log x$, $dB/dL = 2L - 1.1 > 0$ for $L > 0.55$, and $L \mapsto \log x$ is strictly
increasing; $e^{0.55} = 1.7333 < 2$.
\end{proof}

Lemma~\ref{lem:Bmono} is the whole point: $B$ is a function of $p$ \emph{alone}, not of the
pair $(p_n, n)$, so the non-monotone object $T$ never appears in a comparison between two
different indices. In particular this section is independent of (P6$'$), and any downstream
argument that gates the verified range on (P6$'$) is importing a dependency that is not
there.

\subsection{The finite-range theorem}

\begin{theorem}[Unconditional finite-range verification; $\cP$]
\label{thm:range}
Let $X_0 := 60\,184$ and let $X \ge X_0$. Suppose:
\begin{enumerate}
\item[\textup{(H1)}] $p_{n+1}^{\,n} < p_n^{\,n+1}$ for every $n$ with $p_n < X_0$;
\item[\textup{(H2)}] a table is available listing, for every gap value $g$ occurring as
  $g_n = g$ for some $n$ with $p_n \le X$, the first occurrence
  $q(g) := \min\{p_n : g_n = g\}$;
\item[\textup{(H3)}] for every gap value $g$ in that table with $q(g) \ge X_0$, one has
  $g \le B(q(g))$.
\end{enumerate}
Then $\Fb$ holds at every $n$ with $p_n \le X$.
\end{theorem}

\begin{proof}
\idea{A gap that was already safe where it \emph{first} occurred is safe everywhere later,
because the floor $B$ only rises; gaps whose first occurrence is below $X_0$ are all too
small to matter.}
Let $n$ satisfy $p_n \le X$. If $p_n < X_0$, then (H1) and Lemma~\ref{lem:reduction} give
$\Fb$ at $n$. Otherwise $p_n \ge X_0$. Put $g := g_n$ and $q := q(g) \le p_n$.

If $q \ge X_0$: by (H3), $g \le B(q)$; by Lemma~\ref{lem:Bmono} and $q \le p_n$,
$B(q) \le B(p_n)$; by Lemma~\ref{lem:floor}, $B(p_n) < T_n$. Hence $g_n < T_n$, and
Lemma~\ref{lem:reduction} gives $\Fb$ at $n$.

If $q < X_0 \le p_n$: the first occurrence of $g$ lies in the base-case range, so
$g \le G_0 := \max\{g_k : p_k < X_0\} = 72$, attained at $p = 31\,397$
(\S\ref{sec:computation}, item V1, exact integer arithmetic over the $6\,076$ gaps
below $X_0$). Since $B(X_0) = B(60\,184) = 109.008\ldots > 72 \ge g$, Lemmas~\ref{lem:Bmono}
and \ref{lem:floor} give $g \le G_0 < B(X_0) \le B(p_n) < T_n$, so $\Fb$ at $n$.
\end{proof}

\begin{corollary}[Decision procedure; $\cP$]
\label{cor:S}
For a gap value $g$ define
\[
S^\sharp(g) \;:=\; \exp\!\left(\frac{1.1 + \sqrt{1.21 + 4g}}{2}\right),
\]
the unique $p$ with $B(p) = g$, so that $g \le B(p) \iff p \ge S^\sharp(g)$. Then
hypothesis \textup{(H3)} reads $q(g) \ge S^\sharp(g)$, and Theorem~\ref{thm:range} says:
\emph{a gap of size $g$ occurring at a prime $p \ge \max(S^\sharp(g), 60\,184)$ cannot
violate $\Fb$.}
\end{corollary}

\begin{proof}
\idea{Invert the quadratic $B(p) = \log^2 p - 1.1\log p$ in the variable $\log p$ and use
Lemma~\ref{lem:Bmono}.}
Solving $L^2 - 1.1L = g$ for $L > 0$ gives $L = (1.1 + \sqrt{1.21+4g})/2$; $B$ is strictly
increasing (Lemma~\ref{lem:Bmono}), so $g \le B(p)$ iff $\log p$ exceeds that root. The
restatement of (H3) is immediate, and the displayed sentence is the second case analysis of
Theorem~\ref{thm:range} read backwards.
\end{proof}

Note that a gap size which passes is safe \emph{everywhere}, not merely below $X$:
$S^\sharp(g)$ does not depend on $X$. The verification of a range therefore costs one finite
integer check below $60\,184$, plus one comparison per distinct gap value. Nothing else.

\subsection{Independent reproduction of the published constant $1920$}

\begin{proposition}[$\cC$]
\label{prop:1920}
$B(2^{64}) = L(L-1.1)\bigr|_{L = \log 2^{64}} = 1919.1379834975\ldots$, so a gap of at least
$1920$ is required to violate $\Fb$ at a prime just below $2^{64}$. Equivalently, the largest
even $g$ with $S^\sharp(g) \le 2^{64}$ is $g = 1918$:
$S^\sharp(1918) = 1.8209\cdot10^{19} \le 2^{64} = 1.8447\cdot10^{19} < S^\sharp(1920) =
1.8629\cdot10^{19}$.
\end{proposition}

Kourbatov's 2023 endnote states that ``prime gaps of size $g < 1920$ cannot violate''
the conjecture \cite{kourbatov2015verification}; since every prime gap after $g_1 = 1$ is
even, $g < 1920$ means exactly $g \le 1918$. The published integer falls out of
Lemma~\ref{lem:floor} with no tuning.

\begin{hazard}[What the agreement is worth, and what it is not]
\label{haz:1920}
Three caveats, all stated rather than glossed.
(a) The two derivations are not independent in their analytic ingredient: both replace $n$ by
an upper bound on $\pi(p_n)$; they differ in \emph{which} bound (Dusart eq.~(6.6) here, the
Axler-based $L^2 - L - 1.17$ there). What the agreement does establish is that no sign error,
inverted inequality or transcription slip separates the two chains.
(b) Under Kourbatov's sharper constant the same computation yields $1922$, so his published
$g < 1920$ is conservative relative to his own criterion, and the exact coincidence with
$1918$ is partly arithmetic luck at even-gap granularity.
(c) Lemma~\ref{lem:floor} yields the \emph{local} statement at the frontier; the endnote's
phrasing is \emph{global}. Since $S^\sharp(1918) = 1.82\cdot10^{19}$, our bound leaves a gap
of $1918$ free to violate $\Fb$ anywhere below $1.82\cdot10^{19}$. Closing that needs the
first-occurrence table again.
\end{hazard}

\subsection{The table-free window, and its permanent closure}

Can a range be certified with \emph{no} prime enumeration at all? Only just, and only once.

\begin{proposition}[Table-free window; $\cPs$]
\label{prop:window}
$\Fb$ holds at every $n$ with $396\,738 \le p_n \le 777\,600$, with no computational input
beyond two cited explicit estimates.
\end{proposition}

\begin{proof}
\idea{Dusart's short-interval theorem caps $g_n$ by $p/(25L^2)$; Lemma~\ref{lem:floor} floors
$T_n$ by $L^2 - 1.1L$; the two cross exactly once.}
For $p_n \ge 396\,738$, \cite[Prop.~6.8]{dusart2010estimates} supplies a prime in
$(x, x(1 + 1/(25\ln^2 x))]$, hence $g_n \le p_n/(25 L_n^2)$. By Lemma~\ref{lem:floor},
$T_n > L^2 - 1.1L$. So it suffices that
\[
\frac{p}{25L^2} \;\le\; L^2 - 1.1L
\qquad\Longleftrightarrow\qquad
p \;\le\; 25 L^3(L - 1.1), \qquad L = \log p .
\]
Let $h(p) := p - 25(\log p)^3(\log p - 1.1)$. Then $h(396\,738) = -2.35\cdot10^5 < 0$, and
$h'(p) = 1 - (25/p)(4L^3 - 3.3L^2) > 0$ throughout $[4\cdot10^5,\infty)$, so $h$ has a unique
sign change there, at $p^* = 777\,600.744\ldots$ (\S\ref{sec:computation},
item V6). Hence the displayed inequality holds on $[396\,738, p^*]$, so
$g_n < T_n$, so $\Fb$ at $n$ by Lemma~\ref{lem:reduction}.
\end{proof}

\begin{proposition}[The window closes permanently; $\cP$]
\label{prop:closure}
For $p > p^* = 777\,600.744\ldots$, \cite[Prop.~6.8]{dusart2010estimates} does not imply
$\Fb$, and the shortfall grows without bound:
$\dfrac{p/(25L^2)}{L^2 - 1.1L} = \dfrac{p}{25L^3(L-1.1)} \to \infty$.
\end{proposition}

\begin{proof}
\idea{$p$ outgrows every fixed power of $\log p$.}
Immediate from the display in the proof of Proposition~\ref{prop:window}.
\end{proof}

\begin{remark}[The precise measure of what analysis alone can do]
\label{rem:analysis-alone}
A table-free proof of $\Fb$ on $[X_0, X]$ requires an unconditional explicit gap bound of
quality $g_n = O(\log^2 p_n)$ on that range. Every known unconditional gap bound --- Dusart's
$p/(25\log^2 p)$ here, $p^{0.525}$ of \cite{baker2001difference}, and everything between ---
is a \emph{power of $p$}. So reading (B) is dead above $7.776\cdot10^5$ and will stay dead
until unconditional prime-gap technology reaches $\log^2$ scale, which is exactly the strength
that proving $\Fb$ itself requires. \textbf{Corollary, stated so that it is not re-derived:
no amount of analysis will ever eliminate the computational input from a verified-range
claim.} The verified range is irreducibly a computation plus a theorem about the computation.
Theorem~\ref{thm:range} is the theorem; \S\ref{sec:limits} names the computation.
The window $[396\,738, 777\,600]$ lies entirely above $X_0$ and is therefore already covered
by (H1)--(H3) in any real verification; its value is that it is the exact measure of how much
analysis alone can do, and the measure is: one factor of about $2$ in $p$, once, at
$p \approx 10^6$.
\end{remark}

% ==========================================================================
\section{The Riemann-Hypothesis route is closed}
\label{sec:rh}

The slogan ``RH does not help with Firoozbakht'' is common and usually unargued. This section
replaces it with theorems. Five inequivalent readings of ``an RH-conditional bound settles
$\Fb$'' are separated; four are refuted, and the fifth is left undecided with a lower bound
on the strength of what a proof would have to yield.

Throughout this section we use the CMS envelope
\[
B_n := \tfrac{22}{25}\sqrt{p_n}\, L_n ,
\]
for which \cite[\S1.2, Thm.~5]{carneiro2019fourier} gives: RH implies $g_n \le B_n$ for every
$n \ge 3$ (equivalently $p_n > 3$; see \cite[Thm.~1, eq.~(1.4)]{visser2018andrica} for the
same statement with hypothesis $n \ge 3$, $p_n \ge 5$).

We say the envelope \emph{certifies} $\Fb$ at $n$ when $B_n \le T_n$; by
Remark~\ref{rem:strict} this upgrades to the strict inequality $g_n < T_n$ that
Lemma~\ref{lem:reduction} demands.

We also use the following in-run fact, verified for all $n \le 216\,815$ and asymptotically a
consequence of \cite[Thm.~6.9]{dusart2010estimates}: $\{n : T_n \ge L_n^2\} = \{1,\dots,7,10\}$.

\begin{lemma}[$\cP$]
\label{lem:A1}
For every real $x > 0$, $\sqrt{x} > \tfrac{25}{22}\log x$; equivalently
$\tfrac{22}{25}\sqrt{x}\,\log x > (\log x)^2$ for all $x > 1$.
\end{lemma}

\begin{proof}
\idea{One stationary point, and it sits above zero.}
Put $k := 25/22$ and $h(x) := \sqrt{x} - k\log x$ on $(0,\infty)$. Then
$h'(x) = (\sqrt{x} - 2k)/(2x)$, so $h$ has a unique stationary point at
$x^* = (2k)^2 = 625/121 = 5.16528925\ldots$, a minimum, with $h \to +\infty$ at both ends. At
that point
\[
h(x^*) = \tfrac{25}{11} - \tfrac{25}{22}\log\tfrac{625}{121} = 0.40686238\ldots > 0 .
\]
Hence $h > 0$ throughout. Multiplying $\sqrt{x} > k\log x$ by $\log x > 0$ (valid for
$x > 1$) and dividing by $k$ gives the second form.
\end{proof}

\begin{theorem}[The CMS envelope certifies at one index; $\cP$]
\label{thm:A}
Let $S := \{n \ge 3 : B_n \le T_n\}$. Then $S = \{3\}$. That is: the sharpest published
RH-conditional prime-gap bound certifies the Firoozbakht inequality at the single index
$n = 3$ ($p_3 = 5$) and at no other index in the range where the bound is available.
\end{theorem}

\begin{proof}
\idea{Off a set of six indices the bar $T_n$ is below $L_n^2$, and the envelope is above
$L_n^2$ always; the six are checked by hand.}
Let $n \ge 3$. If $n \notin \{3,4,5,6,7,10\}$ then $T_n < L_n^2$, while
Lemma~\ref{lem:A1} at $x = p_n \ge 5$ gives $B_n > L_n^2 > T_n$, so $n \notin S$. The six
remaining indices are checked directly in $60$-digit arithmetic:
\begin{center}
\begin{tabular}{rrrrc}
\toprule
$n$ & $p_n$ & $T_n$ & $B_n$ & $B_n < T_n$? \\
\midrule
3 & 5 & 3.549879733 & 3.166955068 & yes \\
4 & 7 & 4.386035932 & 4.530587009 & no \\
5 & 11 & 6.769336928 & 6.998568638 & no \\
6 & 13 & 6.934281081 & 8.138289656 & no \\
7 & 17 & 8.481637825 & 10.27984133 & no \\
10 & 29 & 11.61044961 & 15.95742984 & no \\
\bottomrule
\end{tabular}
\end{center}
All six comparisons are strict, so $\le$ and $<$ agree on them. Hence $S = \{3\}$.
\end{proof}

\begin{remark}[Read the quantifier]
\label{rem:quantifier}
The restriction $n \ge 3$ in Theorem~\ref{thm:A} is part of the statement and must not be
dropped when the theorem is quoted. The envelope in fact clears the bar at $n = 1$ and
$n = 2$ as well; what excludes them is the \emph{source's} hypothesis $p_n > 3$, not the
arithmetic. The correct informal form is therefore ``at exactly one index in the range where
the CMS bound is available'', never ``at no other index whatsoever''.
\end{remark}

\begin{remark}
The single certified index is worthless: $n = 3$ means $\Fb$ at $p = 5$, i.e.\
$7^3 = 343 < 5^4 = 625$, decidable by hand, and $n = 3$ falls below the threshold $k > 9$
that both Kourbatov criteria carry by hypothesis. Moreover RH is not even load-bearing where
the envelope is smallest: by \cite[\S7]{visser2018andrica} the CMS inequality holds
\emph{unconditionally} for all $p < 1.836\cdot10^{19}$, essentially the whole verified range
of Theorem~\ref{thm:verified} --- and what is already known still certifies $\Fb$ at no index
above $n = 3$.
\end{remark}

\begin{theorem}[No power-type envelope suffices; $\cP$]
\label{thm:Bpow}
Let $C > 0$, $\theta > 0$, $A \in \R$ and $E(x) := C x^\theta (\log x)^A$. Then
$\#\{n : E(p_n) \le T_n\} < \infty$; indeed $E(p_n)/T_n \to \infty$.
\end{theorem}

\begin{proof}
\idea{In the variable $u = \log x$ the envelope is exponential and the bar is quadratic.}
Since $T_n < L_n^2$ for all $n \notin \{1,\dots,7,10\}$, it suffices that
$E(x)/(\log x)^2 \to \infty$. Write $u := \log x$; then
$E(x)/(\log x)^2 = C e^{\theta u} u^{A-2} \to \infty$ as $u \to \infty$, because $e^{\theta u}$
with $\theta > 0$ dominates any fixed power of $u$. Hence
$E(p_n)/T_n > E(p_n)/L_n^2 \to \infty$, and a divergent ratio is $\le 1$ only finitely often.
\end{proof}

Every prime-gap upper bound currently available lives in this class:
$p_n^{0.525}$ unconditionally \cite{baker2001difference}; $O(\sqrt{p_n}\log p_n)$ under RH
(Cram\'er 1919, as reported in \cite[\S2, Thm.~4]{visser2018andrica}); $4\sqrt{p_n}\log p_n$
under RH (Goldston 1982, \cite[\S2, Thm.~5]{visser2018andrica}); the explicit CMS bound and
its asymptotic refinement $\limsup \le 1/C^+(B) < 21/25$
\cite[\S1.2, Cor.~4]{carneiro2019fourier}. \textbf{The column that matters is $\theta$:} every
entry has $\theta > 0$, while the bar \eqref{eq:sufficient} has $\theta = 0$. Improving the
exponent from $0.525$ to $0.5$ --- exactly what RH buys --- moves nothing across the line,
because the line is at $\theta = 0$.

\begin{theorem}[The critical constant is $2/e$; $\cP$]
\label{thm:Ccrit}
For $C > 0$ set $E_C(x) := C\sqrt{x}\log x$. Then
$\{x > 1 : E_C(x) \le (\log x)^2\} \ne \emptyset$ if and only if $C \le 2/e = 0.735758882\ldots$,
and when non-empty this set is the \emph{bounded} interval with endpoints
$\exp(-2W_0(-C/2))$ and $\exp(-2W_{-1}(-C/2))$, where $W_0, W_{-1}$ are the two real branches
of the Lambert $W$ function. Consequently, for every $C > 0$, $E_C$ certifies $\Fb$ at only
finitely many indices.
\end{theorem}

\begin{proof}
\idea{The condition is $C \le \log x/\sqrt{x}$, and that function is unimodal with maximum
$2/e$ at $x = e^2$; solving the equality is a Lambert-$W$ substitution.}
$E_C(x) \le (\log x)^2 \iff C \le \varphi(x) := \log(x)/\sqrt{x}$ for $x > 1$. Now
$\varphi'(x) = (2 - \log x)/(2x^{3/2})$, so $\varphi$ increases on $(0,e^2)$, decreases on
$(e^2,\infty)$, and attains its maximum $\varphi(e^2) = 2/e$ at $x = e^2 = 7.389056\ldots$.
Hence the superlevel set $\{x > 1 : \varphi(x) \ge C\}$ is non-empty iff $C \le 2/e$ and, by
unimodality with $\varphi \to 0$ at both ends, is a compact interval. Solving
$C\sqrt{x} = \log x$ with $t = \log x$ gives $te^{-t/2} = C$, i.e.\
$(-t/2)e^{-t/2} = -C/2$, i.e.\ $-t/2 = W(-C/2)$, whence the stated endpoints; two real
branches exist iff $-C/2 \ge -1/e$, i.e.\ $C \le 2/e$. Finally the set is bounded, so it
contains finitely many primes, and away from it $E_C(x) > (\log x)^2 > T_n$ off the exception
set $\{1,\dots,7,10\}$.
\end{proof}

Both published constants, $22/25 = 0.88$ and $21/25 = 0.84$, are \emph{above} the critical
$2/e$, so both give the empty set. Constants far below it buy a bounded initial segment and
then stop: $C = 0.1$ covers $p \in (1.111, 8099)$, and $C = 0.01$ covers
$p \in (1.010, 2\,122\,265)$. To certify $\Fb$ merely up to the \emph{already verified}
frontier one would need $C \le (L^2 - L - 1)/(\sqrt{p}L) = 1.009\cdot10^{-8}$ at $p = 2^{64}$
--- not a constant improvement but the abandonment of the $\sqrt{p}$ scale.

\begin{corollary}[Even $\limsup = 0$ does not help; $\cP$]
\label{cor:limsup0}
Suppose $g_n = o(\sqrt{p_n}\log p_n)$. This certifies $\Fb$ at no index.
\end{corollary}

\begin{proof}
\idea{$o(\sqrt{p}\log p)$ does not even imply $g_n < \log^2 p_n$, so it cannot imply
$g_n < T_n$.}
The hypothesis is compatible with gaps far above the threshold: a sequence with
$g_n = \sqrt{p_n}L_n/\log\log p_n$ satisfies $g_n/(\sqrt{p_n}L_n) \to 0$ while
$g_n/L_n^2 = \sqrt{p_n}/(L_n \log\log p_n) \to \infty$. Killing the constant does not change
the scale, and by Theorem~\ref{thm:Bpow} the scale is the whole obstruction.
\end{proof}

\begin{hazard}[Provenance of the hypothesis of Corollary~\ref{cor:limsup0}]
That $\limsup = 0$ holds under RH together with Montgomery's pair correlation conjecture is
\emph{reported}, not proved, by \cite[\S1.2, after eq.~(1.14)]{carneiro2019fourier}, who
attribute it to three works not opened in this study. It is therefore second-hand and is
attributed as such here. Corollary~\ref{cor:limsup0} does not depend on the statement being
true --- it says that \emph{even if} it holds, nothing is certified --- so the weak tier does
not propagate into the verdict.
\end{hazard}

\subsection{The material implication $\mathrm{RH} \Rightarrow \Fb$}

Theorems~\ref{thm:A}--\ref{thm:Ccrit} refute the \emph{route}. They say nothing about the
\emph{proposition}. Conflating the two is the failure mode this section most needs to avoid,
so we state plainly: to refute $\mathrm{RH} \Rightarrow \Fb$ one must establish
$\mathrm{RH} \wedge \neg\Fb$, i.e.\ prove the Riemann Hypothesis \emph{and} exhibit a
counterexample to $\Fb$; to prove it one must derive $\Fb$ from RH. Neither is available and
no computation bears on either. \textbf{The implication is undecided here and we do not
pretend otherwise.} What can be proved is a lower bound on the strength of what such a proof
would yield.

\begin{theorem}[$\cP$]
\label{thm:D}
Suppose $\mathrm{RH} \Rightarrow \Fb$ were proved. Then, as an immediate corollary, one would
have proved
\[
\mathrm{RH} \;\Longrightarrow\; \bigl(\forall k > 9 : g_k < L_k^2 - L_k - 1\bigr),
\]
an RH-conditional Cram\'er-type prime-gap bound at $\log^2$ scale, uniform in $k$, with
leading constant $1$.
\end{theorem}

\begin{proof}
\idea{Compose with the unconditional necessary condition \eqref{eq:necessary}.}
Compose the hypothesised implication with \cite[\S2, Thm.~1]{kourbatov2015bounds}, which is
the unconditional implication \eqref{eq:necessary}.
\end{proof}

\begin{corollary}[The size of the required advance; $\cP$]
\label{cor:D1}
The ratio of the best published RH-conditional bound to the bound Theorem~\ref{thm:D} would
deliver is
\[
\frac{\tfrac{22}{25}\sqrt{p_n}\,L_n}{L_n^2 - L_n - 1} \;\longrightarrow\; \infty ,
\qquad\text{the ratio growing like } \frac{22}{25}\cdot\frac{\sqrt{p_n}}{L_n} ,
\]
with value $8.72\cdot10^7$ at $p = 2^{64}$ and unbounded thereafter.
\end{corollary}

\begin{hazard}[No difficulty ordering is claimed]
\label{haz:noorder}
Corollary~\ref{cor:D1} is a \emph{distance between two statements}, not a measure of proof
effort. Nothing here says such a proof must first pass through a sharpened $\sqrt{p}$-scale
bound, and nothing here orders $\mathrm{RH}\Rightarrow\Fb$ against any other open problem by
difficulty. In particular this paper does \emph{not} assert that $\Fb$ is harder than RH.
\end{hazard}

\begin{corollary}[The circularity charge, made quantitative; $\cPs$]
\label{cor:D2}
Let $H$ be any hypothesis of gap-bound shape with $H \Rightarrow \Fb$. Composing with
\eqref{eq:necessary} gives $H \Rightarrow (\forall k > 9 : g_k < L_k^2 - L_k - 1)$. So every
sufficient hypothesis is at least as strong as the necessary condition; and by
\eqref{eq:sufficient} a hypothesis only an additive $0.17$ stronger than that bound already
suffices. Hence a candidate gap-bound hypothesis can sit only in the band between
$L^2 - L - 1$ and $L^2 - L - 1.17$: outside it, a candidate is either too weak to give $\Fb$
or is itself a sufficient condition, i.e.\ $\Fb$ with a constant changed.
\end{corollary}

\begin{remark}
Corollary~\ref{cor:D2} does \emph{not} say that every sufficient $H$ is within $0.17$ of
$\Fb$ --- plainly false, since $g_k < L^2 - L - 5$ is sufficient and far stronger. The
sandwich pins $\Fb$, not the class of hypotheses that imply it. Corollary~\ref{cor:D2} uses
\eqref{eq:sufficient}, hence Axler (Caveat~\ref{haz:axler}); the inequality \emph{direction}
is robust to the constant --- any $b$ with $1 < b < \infty$ gives the same conclusion with
$b - 1$ in place of $0.17$ --- but the numeral $0.17$ inherits that dependency. The
corollary is also stated for hypotheses of gap-bound shape only.
\end{remark}

\subsection{Cram\'er's hypothesis does not entail $\Fb$ over integer sequences}

The last escape is: assume Cram\'er's conjecture instead of RH. Write
\begin{equation}
\label{eq:Cr}
(\mathrm{Cr}) \qquad \limsupn\, g_n / L_n^2 \;\le\; 1 .
\end{equation}

\begin{lemma}[$\cP$]
\label{lem:E1}
For every $\delta > 0$ there are infinitely many $n$ with $g_n < (1+\delta)L_n$.
\end{lemma}

\begin{proof}
\idea{If all gaps were eventually $\ge (1+\delta)\log p_n$, the primes would be sparser than
the prime number theorem allows.}
Suppose not: $g_n \ge (1+\delta)L_n$ for all $n \ge N_0$. Summing,
$p_N \ge (1+\delta)\sum_{n=N_0}^{N-1} L_n$. By the prime number theorem $p_n \sim n\log n$,
so $L_n \sim \log n$ and $\sum_{n \le N} L_n \sim N\log N$. Hence
$p_N \ge (1 + \delta + o(1))N\log N$, contradicting $p_N \sim N\log N$ for large $N$.
\end{proof}

\begin{theorem}[Counter-model; $\cP$]
\label{thm:E}
Assume $(\mathrm{Cr})$. Then there is a strictly increasing sequence $(q_n)_{n\ge1}$ of
positive integers with
\begin{enumerate}
\item $q_n = p_n + O\bigl((\log n)^2 \log\log n\bigr)$;
\item $\limsupn (q_{n+1}-q_n)/(\log q_n)^2 = 1$, so $q$ satisfies $(\mathrm{Cr})$;
\item $q_{n+1}^{1/(n+1)} \ge q_n^{1/n}$ for infinitely many $n$, so $q$ violates the
  Firoozbakht inequality infinitely often.
\end{enumerate}
Hence $(\mathrm{Cr})$ does not entail $\Fb$ in the theory of strictly increasing sequences of
positive integers.
\end{theorem}

\begin{proof}
\idea{Take the primes and, at a sparse sequence of indices where the gap happens to be
small, inflate that one gap up to exactly $\lceil \log^2 \rceil$; the inflations are so rare
and so small that they do not disturb the $\limsup$ or the density, but each one breaks the
Firoozbakht inequality.}
By Lemma~\ref{lem:E1} with $\delta = 1$, for each $k \ge 1$ let
$n_k := \min\{n \ge 2^{2^k} : g_n \le 2L_n\}$, which exists, and $n_k \to \infty$. Define
recursively
\[
J_k := \bigl\lceil (\log q_{n_k})^2 \bigr\rceil - g_{n_k},
\qquad
q_n := p_n + \sum_{k\,:\,n_k < n} J_k ;
\]
the recursion is well founded because $q_{n_k}$ depends only on $J_1,\dots,J_{k-1}$.

\emph{Positivity and monotonicity.} $g_{n_k} \le 2L_{n_k}$ while
$\lceil (\log q_{n_k})^2\rceil \ge L_{n_k}^2$, and $L^2 > 2L$ for $L > 2$, i.e.\ $p > e^2$.
Start the construction at the least $k_0$ with $p_{n_{k_0}} \ge 11$ and discard $k < k_0$;
this is legitimate because the conclusion concerns infinitely many indices, and it makes
every $J_k \ge 0$. Then $q$ is $p$ plus a nondecreasing step function, hence strictly
increasing.

\emph{Item 1.} $q_n - p_n = \sum_{k : n_k < n} J_k$ with
$J_k \le \lceil (\log q_{n_k})^2\rceil$. Since $n_k \ge 2^{2^k}$, the number of terms with
$n_k < n$ is at most $\log_2\log_2 n + O(1)$, and each $J_k \le (\log q_n)^2 + 1 = O((\log n)^2)$.
Hence $q_n - p_n = O((\log n)^2\log\log n)$.

\emph{Item 2.} At $n = n_k$, $q_{n_k+1} - q_{n_k} = g_{n_k} + J_k = \lceil(\log q_{n_k})^2\rceil$
by construction, so the ratio is $\lceil(\log q_{n_k})^2\rceil/(\log q_{n_k})^2 \to 1$. At
$n \notin \{n_k\}$, $q_{n+1}-q_n = g_n$ and $\log q_n \ge \log p_n$, so the ratio is at most
$g_n/L_n^2$, whose $\limsup$ is $\le 1$ by $(\mathrm{Cr})$. Hence the $\limsup$ is exactly $1$.

\emph{Item 3.} Write $T^{(q)}_n := q_n(q_n^{1/n}-1)$; the derivation of
Lemma~\ref{lem:reduction} is purely formal and applies to any strictly increasing positive
sequence, so the Firoozbakht inequality fails at $n$ iff $q_{n+1}-q_n \ge T^{(q)}_n$. The
index $n$ is unchanged by the construction, and $q_n = p_n(1+\varepsilon_n)$ with
$\varepsilon_n = O((\log n)^2\log\log n/p_n)$, which tends to $0$ faster than any negative
power of $L_n$. Since $\partial T_n/\partial p_n = O(L_n^2/p_n)$, the perturbation moves the
bar by $O(\varepsilon_n L_n^2) = o(1)$; with $T_n = L_n^2 - L_n - 1 + o(1)$
\cite[\S5, Thm.~5]{kourbatov2015bounds} this gives $T^{(q)}_n < (\log q_n)^2$ for large $n$.
At $n = n_k$ the gap is $\lceil(\log q_{n_k})^2\rceil \ge (\log q_{n_k})^2 > T^{(q)}_{n_k}$.
So the inequality fails at every sufficiently large $n_k$.
\end{proof}

\begin{remark}[What Theorem~\ref{thm:E} does and does not close]
\label{rem:E-scope}
It closes the reading ``Cram\'er's $\limsup$ hypothesis entails $\Fb$'' \emph{as a
counter-model result}. It does \emph{not} prove $(\mathrm{Cr}) \not\Rightarrow \Fb$ as a
material implication about the primes: $p$ is one fixed sequence, and if both
$(\mathrm{Cr})$ and $\Fb$ happen to be true the implication is vacuously true.
Theorem~\ref{thm:E} quantifies over sequences, on purpose. The stronger informal reading ---
``no derivation using only growth and distribution can succeed'' --- is a \emph{gloss}, not a
corollary: the drift in item~1 is small enough that no unconditional $\pi(x)$ estimate in our
toolbox separates $q$ from the primes (the induced displacement of the counting function is
$O(\log x\log\log x)$, against a Dusart bracket of width $0.2762\,x/\log^2 x$), so a
derivation appealing only to $(\mathrm{Cr})$, the prime number theorem and such estimates
would have to prove $\Fb$ for $q$ as well. That argument is only as good as its unformalised
premise about which appeals a derivation makes, and the class it names is contingent on what
this study fetched. \textbf{Quote the theorem, not the gloss.}
\end{remark}

Structurally, Theorem~\ref{thm:E} explains why the additive $0.17$ of
Corollary~\ref{cor:D2} is not a technicality: $(\mathrm{Cr})$ controls a $\limsup$, whereas
$\Fb$ needs a uniform bound at \emph{every} index with the second-order term pinned. The
distance between ``asymptotically at most $1$'' and ``below $L^2 - L - 1$ always'' is exactly
the room the construction lives in.

% ==========================================================================
\section{The smooth model}
\label{sec:smooth}

It is worth recording why the difficulty is entirely in the fluctuation.

\begin{proposition}[$\cP$]
\label{prop:smooth}
The function $x \mapsto (x\log x)^{1/x}$ is strictly decreasing on $x \ge 5$.
\end{proposition}

\begin{proof}
\idea{Take logarithms; the derivative of $\log(x\log x)/x$ has the sign of
$1 + 1/\log x - \log(x\log x)$, which is negative once $x$ is large enough, and $x = 5$ is
large enough.}
Write $\psi(x) := \log(x\log x)/x$. Then
$x^2\psi'(x) = \bigl(1 + 1/\log x\bigr) - \log(x\log x)$. The bracket
$1 + 1/\log x$ decreases and $\log(x\log x)$ increases, so $\psi'$ has at most one sign
change; evaluating, the bracket at $x=5$ is $1.6213$ while $\log(5\log 5) = 2.0853$, so
$\psi'(5) < 0$, and $\psi' < 0$ for all $x \ge 5$. Hence $\psi$, and therefore
$(x\log x)^{1/x} = e^{\psi(x)}$, is strictly decreasing on $[5,\infty)$.
\end{proof}

Since $p_n \sim n\log n$, Proposition~\ref{prop:smooth} says that the \emph{smooth model} of
Firoozbakht's conjecture is true and elementary. The entire difficulty of $\Fb$ localises in
the fluctuation of $p_n$ around $n\log n$. The safe stated range is $x \ge 5$, not $x \ge 4$:
the relevant bracket changes sign just above $4$ ($+0.0084$ at $x=4$, $-0.0193$ at $x = 4.05$).
Proposition~\ref{prop:smooth} was \emph{not} formalised in Lean in this work
(\S\ref{sec:lean}).

% ==========================================================================
\section{Computational evidence}
\label{sec:computation}

\subsection{Design}

Two independent legs, with independent code paths, swept the primes exhaustively to $10^{11}$
--- $4\,118\,054\,812$ consecutive prime pairs, no pruning, no sampling, no assumed lemma.
The protocol is two-tier and the tiering is the point:

\begin{itemize}
\item \textbf{Screen.} \texttt{float64} over all pairs, on $\rho_n = g_n/T_n$. This statistic
  has $O(1)$ operands and so does not carry the catastrophic cancellation that afflicts the
  ratio statistic of \S\ref{sec:artefact}.
\item \textbf{Escalation.} Every pair with $\rho_n \ge 0.90$ is re-decided in arbitrary
  precision with an explicit error budget, and the certified verdict function \emph{raises}
  rather than returns when the margin lands inside that budget.
\item \textbf{Calibration.} The high-precision path is checked against the exact integer
  criterion $p_{n+1}^{\,n} < p_n^{\,n+1}$ for every $n \le 5000$ in one leg and $n \le 4000$
  in the other. Zero disagreements.
\end{itemize}

The escalation discipline deserves emphasis because it inverts the usual default:
\emph{the silent failure here is in the verification direction.} A probe that only breaks on
a detected violation reports ``no counterexample'' out of numerical noise. This one cannot.

\subsection{Results}

\begin{center}
\begin{tabular}{@{}lr@{}}
\toprule
Quantity & Value \\
\midrule
consecutive prime pairs checked ($p_n < 10^{11}$) & $4\,118\,054\,812$ \\
violations of $\Fb$ & \textbf{0} \\
$\max_{n \ge 10}\rho_n$ & $\mathbf{0.8317570}$ at $p_n = 25\,056\,082\,087$, $g = 456$ \\
$\max_{n \ge 10} g_n/L_n^2$ & $0.7953487$, same prime \\
largest gap in range & $464$ at $p = 42\,652\,618\,343$ \\
maximal-gap records in range & $40$ \\
sieve validation & $\pi(10^9)$, $\pi(10^{11})$ match standard values \\
\bottomrule
\end{tabular}
\end{center}

Only two pairs in the whole range reach $\rho \ge 0.9$, and both are at $n < 5$
($\rho_2 = 0.9107$, $\rho_4 = 0.9120$), which is why the record is quoted for $n \ge 10$.

\subsection{Verification of the lemmas of \S\ref{sec:monotone}--\S\ref{sec:range}}

The following items are cited by number from the proofs above. All were recomputed
independently of the code paths that produced the sweep.

\begin{enumerate}
\item[V1] $\max\{g_k : p_k < 60\,184\} = 72$, attained at $p = 31\,397$, by
  exact integer arithmetic over the $6\,076$ gaps below $60\,184$; and
  $B(60\,184) = 109.008\ldots > 72$. Also: $p_{n+1}^{\,n} < p_n^{\,n+1}$ for every one of
  those $6\,076$ indices in exact integer arithmetic --- this is hypothesis~(H1) of
  Theorem~\ref{thm:range}, discharged.
\item[V2] $S(29) = 6.8013853766\ldots > 6 = \max\{g_j : j \le 9\}$ ---
  hypothesis~(ii) of Theorem~\ref{thm:B}.
\item[V3] Lemma~\ref{lem:floor} holds at every prime in $[60\,184, 10^{11}]$; the tightest
  slack is $T_n - B(p_n) = +0.0798914728\ldots$, attained at the very bottom of the range
  ($p = 155\,893$).
\item[V4] Lemma~\ref{lem:dusart} holds at every prime in $[60\,184, 2\cdot10^7]$; zero
  failures.
\item[V5] $B(2^{64}) = 1919.1379834975328$ (Proposition~\ref{prop:1920}), and the largest
  even $g$ with $S^\sharp(g) \le 2^{64}$ is $1918$.
\item[V6] The sign change of $h(p) = p - 25(\log p)^3(\log p - 1.1)$ is at
  $p^* = 777\,600.744\ldots$ (Proposition~\ref{prop:window}).
\item[V7] Table lookup for the whole in-run range: of the $209$ gap sizes occurring below
  $10^{11}$, $55$ are excluded outright by $S^\sharp(g) < 60\,184$ and $154$ need the table
  test; \textbf{none is unsettled}. The minimum safety factor $q(g)/S^\sharp(g)$ is
  $5.337$ at $g = 112$, the maximum $203.9$ at $g = 334$.
\item[V8] At the published record Cram\'er--Shanks--Granville point
  $p = 1\,693\,182\,318\,746\,371$, $g = 1132$ \cite{oeisA111943}: $B(p) = 1191.009$, so the
  record gap is safe with a margin of $4.955\%$.
\item[V9] The certificate that scales: splitting $[60\,184, 10^{11}]$ into $22$ geometric
  windows $[a,b)$ and checking $\max\{g_n : p_n < b\} < B(a)$ certifies every window,
  consuming \emph{only} the $40$ maximal-gap records --- no sieve, no $\pi(x)$ computation,
  no per-prime check. This is the shape in which a frontier like $2^{64}$ is actually
  certified.
\end{enumerate}

\begin{remark}[Stability of the table test across four decades]
The minimum safety factor $q(g)/S^\sharp(g)$ is $5.3371$ at $g = 112$ at every one of
$X = 10^7, 10^8, 10^9, 10^{10}, 10^{11}$. Four decades of new data and $134$ new gap sizes
produced nothing tighter than a case already visible in the first decade ($g = 112$, first
occurring at $p = 370\,261$). Route (b) neither tightens nor loosens as the range extends.
\end{remark}

\subsection{Three refutations of claims about the evidence}
\label{sec:artefact}

Computation refutes nothing about $\Fb$. It did refute three claims about the evidence, and
recording them is the point of this subsection.

\paragraph{(1) The ``$0.9999984$ near-miss'' is an artefact.}
The statistic $\max_n\, n\log p_{n+1}/((n+1)\log p_n)$ reads $0.9999984$ below $3\cdot10^6$
and has been read as evidence that $\Fb$ nearly fails. It is not. In a synthetic universe
where every gap equals $2$ (so that $\Fb$ holds by a mile) the same statistic reads
$0.99999991$ at $n = 10^7$, while $\rho_n$ there is $0.059$ --- an order of magnitude below
the bar. The statistic measures $1/n$. Only $\rho$ is diagnostic:

\begin{center}
\begin{tabular}{@{}lccccc@{}}
\toprule
$X$ & $10^7$ & $10^8$ & $10^9$ & $10^{10}$ & $10^{11}$ \\
\midrule
$\max$ ratio statistic & 0.99999937 & 0.99999994 & 0.999999993 & 0.9999999992 & 0.99999999992 \\
$\max_{n\ge10}\rho_n$ & 0.7605 & 0.7896 & 0.7896 & 0.7896 & 0.8318 \\
\bottomrule
\end{tabular}
\end{center}

\paragraph{(2) ``The tightest cases sit at record gaps'' does not survive the range.}
The observation that the tightest $\rho$ cases occur at record gaps was a print truncation of
the first six entries. At $10^{11}$: $22$ of the top $40$, and the fourth-tightest case among
four billion pairs is \emph{not} at a record index ($g = 444$ at $p = 36\,172\,730\,063$,
$\rho = 0.78502$). Read precisely: if the pruning rule holds, a record-index search still
misses no \emph{counterexample}; what it misses is \emph{near-misses}. So the observation
cannot be used as evidence for the pruning rule that produced it.

\paragraph{(3) More sieve is not more evidence.}
Two decades of extra sieving ($10^8 \to 10^{10}$) produced no new near-miss: the record
$\rho$ sat at $0.7896$ ($g = 210$, $p = 20\,831\,323$) throughout. It moved exactly once, at
$10^{11}$, and it moved \emph{at the next maximal gap}. Sieving buys the next maximal gap and
nothing in between. The record that would matter --- $\rho \approx 0.948$ at
$p \approx 1.693\cdot10^{15}$ \cite{oeisA111943} --- is $4.2$ decades above anything reached
here. The compute that would matter goes into gap \emph{tables}, not prime \emph{counts}.

\subsection{A counter-model for first-failure maximality}
\label{sec:ffm-counter}

\begin{proposition}[$\cC$]
\label{prop:ffm}
There is an explicit strictly increasing sequence of positive integers whose first
Firoozbakht failure occurs at a non-record gap.
\end{proposition}

\begin{proof}
\idea{Take the real primes up to $1327$, then a large jump, then a long dense run of gaps of
$2$, then a jump smaller than the first one; the dense run drives the index up so far that
the later, smaller gap is the one that breaks the inequality.}
Take the primes up to $1327$, then a jump of $40$, then $100$ jumps of $2$, then a jump of
$38$. The first Firoozbakht failure of this sequence is at index $323$ on the gap of $38$,
which is not a record (the earlier jump of $40$ is larger). Verified in exact integer
arithmetic, so rounding cannot be blamed.
\end{proof}

The consequence is sharp: \textbf{any proof of first-failure maximality that does not consume
an arithmetic density input is wrong}, because the statement is false for sequences satisfying
everything except primality. The good news is that a weak input suffices for this mechanism:
the dense run the counter-model needs exceeds the unconditional Montgomery--Vaughan cap
$2y/\log y$ by a factor that \emph{grows} with scale --- $1.41$ at $p \approx 10^3$, $4.18$ at
$4\cdot10^8$, $9.13$ at $10^{18}$ --- so the counter-model is excluded unconditionally, and
increasingly comfortably. (The Montgomery--Vaughan bound is not covered by a source-ledger row
in this work; cf.\ Caveat~\ref{haz:bt}.)

A second, independent route is also closed: bounding $T_m$ above and $T_n$ below with
explicit $\pi(x)$ estimates \emph{at the two endpoints separately} demands record gaps of
size at least $5.3\cdot10^4$ at $p = 10^6$ (Dusart) or $3.7\cdot10^3$ (Axler), against actual
records of size $\approx \log^2 p \approx 176$; and past $\approx 10^{10}$ (Dusart) the
criterion is unsatisfiable at \emph{any} gap size. The bound spread dwarfs the quantity being
resolved, and no sharpening of constants rescues the route.

% ==========================================================================
\section{Formalisation in Lean~4}
\label{sec:lean}

\subsection{What is machine-checked}

The reduction chain
\[
\texttt{Conjecture} \;\leftrightarrow\; \texttt{ConjectureReal} \;\leftrightarrow\;
\bigl(\forall n \ge 1,\ g_n < T_n\bigr)
\]
--- i.e.\ Lemma~\ref{lem:reduction} --- is formalised in Lean~4 against Mathlib and checked
by the kernel. Toolchain: \texttt{leanprover/lean4:v4.29.0}; Mathlib tag \texttt{v4.29.0},
revision \texttt{8a178386\allowbreak ffc0f5fe\allowbreak f0b77738\allowbreak bb5449d5\allowbreak 0efeea95}.

\begin{center}
\small
\begin{tabular}{@{}p{0.46\textwidth} p{0.44\textwidth}@{}}
\toprule
Gate & Result \\
\midrule
\texttt{lake build} exit code & \textbf{0} (1984 jobs) \\
build warnings & \textbf{1} --- the declared open target, and nothing else \\
declarations scanned by the axiom audit & \textbf{60} (exhaustive, list-free) \\
declarations depending on \texttt{sorryAx} & \textbf{1} --- \texttt{Firoozbakht.firoozbakht} \\
live \texttt{sorry} tokens in sources & \textbf{1} \\
stray \texttt{axiom}, \texttt{native\_decide}, \texttt{@[implemented\_by]} or \texttt{unsafe} & \textbf{none} (grep-clean) \\
\bottomrule
\end{tabular}
\end{center}

The single \texttt{sorryAx} dependent is Firoozbakht's conjecture itself, which was never
attempted. \emph{An attack on an open $\Pi_1$ statement whose \texttt{sorry} count reaches
zero would be reporting a fabrication, not a result.} Four of the skeleton's original five
\texttt{sorry}s --- the equivalence-chain lemmas --- were discharged with real proof terms,
and the signatures were mechanically diffed against the skeleton anchor with zero changes,
which is what forecloses the failure mode ``closed a \texttt{sorry} by weakening the
statement''. The equivalences were also shown non-vacuous: $g_1 = 1$ and $T_1 = 2^2 - 2$ are
pinned to numerals.

The audit itself was made exhaustive rather than trusted: a hand-maintained list of
declarations to \texttt{\#print axioms} is a hazard dressed as a check, since a \texttt{sorry}
in a forgotten declaration is invisible. The list-free audit walks the namespace. The checker
was then tested against a corpus of $27$ adversarial statements, all false or ill-formed; all
$27$ behaved as specified. One corpus entry --- a true-but-differently-meant statement
produced by an $\N \to \R$ coercion --- is caught by \emph{no} gate in this study and is named
as such rather than omitted.

\subsection{What is not formalised, and why}

\begin{itemize}
\item \textbf{The smooth model} (Proposition~\ref{prop:smooth}). Fully within Mathlib's
  reach; simply not done.
\item \textbf{The $\limsup$ corollary.} Needs effective $\pi(x)$ bounds that are not assumed
  present in Mathlib.
\item \textbf{Any verified range past $n \le 4$.} This is a hard limit and deserves naming.
  \texttt{Nat.nth} is \emph{noncomputable} \cite{mathlibNatNth} --- not merely slow --- so
  \texttt{decide} cannot produce $p_n$ at all; and Mathlib's prime-specific \texttt{nth} API
  is exactly five \texttt{@[simp]} base lemmas, $\texttt{nth Prime } 0 = 2$ through
  $\texttt{nth Prime } 4 = 11$ \cite{mathlibNatPrimeNth}. Extending requires building
  \texttt{Nat.count}$\leftrightarrow$\texttt{Nat.nth} bridging machinery, plus prime literals
  with \texttt{Nat.Prime} certificates and a ``no prime strictly between'' lemma for each
  consecutive pair. Reporting a larger verified $N$ without that work would be a fabrication.
\item \textbf{Dusart's estimate} (Lemma~\ref{lem:dusart}) is not in Mathlib and would have to
  enter any formalisation of \S\ref{sec:range} as a hypothesis or an axiom --- and be
  declared as such in any paper quoting the result.
\end{itemize}

\begin{hazard}[Indexing]
Mathlib's \texttt{nth} is $0$-indexed ($\texttt{nth Prime } 0 = 2$) while the conjecture as
posed here is $1$-indexed ($p_1 = 2$). Every index threshold in the literature --- $k > 9$,
$n \ge 10$, $n \ge 3645$, $n \ge 5$ --- is stated in the $1$-indexed convention. The
off-by-one must be fixed once, in the statement file; it is the single most likely source of a
silent error in a formalisation of this material.
\end{hazard}

% ==========================================================================
\section{The tension that will not resolve itself}
\label{sec:tension}

The strongest reason to believe $\Fb$ is \emph{false} is not a computation and not a theorem.
Cram\'er's model, in Granville's corrected form, predicts
$\limsup g_n/\log^2 p_n \ge 2e^{-\gamma} \approx 1.1229 > 1$
\cite[preprint p.~12]{granville1995cramer}, which is incompatible with $\Fb$ via
\eqref{eq:necessary}. This is $\cH$: a heuristic, not a test. Granville's own sign is
$\gtrapprox$ --- ``suggests'', not ``proves''.

Symmetrically, the strongest reasons to believe $\Fb$ is \emph{true} are the verified range
(Theorem~\ref{thm:verified}) and the fact that the record Cram\'er--Shanks--Granville ratio
$0.9206$, standing since 1999 \cite{oeisA111943}, sits comfortably below the
Firoozbakht-implied ceiling of $1 - 1/\log p \approx 0.9715$ at that size --- about $5\%$ of
headroom. Neither is a proof and neither is close to one.

\begin{hazard}[Cram\'er did not conjecture $\limsup = 1$ about the primes]
He \emph{proved} it about his urn model, with probability~1
\cite[p.~27]{cramer1936order}, and \emph{suggested} the analogue for the primes
\cite[p.~24, eq.~(4)]{cramer1936order}. Granville frames it the same way. Any sentence of the
form ``Cram\'er conjectured that $\limsup g_n/\log^2 p_n = 1$'' is a compression the source
does not literally support.
\end{hazard}

% ==========================================================================
\section{What remains open, and what this paper does not establish}
\label{sec:limits}

This section is long on purpose. A paper reporting an attack that failed to settle its target
is exactly the paper that gets over-quoted downstream.

\subsection{$\Fb$ itself, and why it is not close}

The load-bearing obstruction is unchanged by anything above: \textbf{any proof of $\Fb$
yields $g_n = O(\log^2 p_n)$ unconditionally}, via \eqref{eq:necessary}. The best known
unconditional gap bound is $g_n = O\bigl(p_n^{0.525}\bigr)$ \cite{baker2001difference}; under RH it
improves only to $\approx\sqrt{p_n}\log p_n$ \cite{carneiro2019fourier}. Both are powers of
$p_n$; $\Fb$ needs polylogarithmic. Any proposal of a ``direct proof'' must say how it clears
this, or it is proposing something already beyond the field.

Compounding it: \textbf{there is no induction mechanism.} $g_n$ is not constrained by
$g_1,\dots,g_{n-1}$; knowing $\Fb$ up to $n$ gives no leverage at $n+1$. Any proposed
inductive proof must first supply the missing mechanism.

On the refutation side, the best large-gap results \cite{ford2016large} reach an iterated-log
factor above $\log n$ but remain a full power of $\log$ below what a counterexample needs.
Explicit constructions fail for a second, independent reason: they place the gap at an
\emph{unspecified} location, while $\Fb$ needs the gap and the count $\pi(p_n) = n$ at the
\emph{same} point. Correspondingly, although $\neg\Fb$ is $\Sigma_1$ and finitely certifiable
in principle, the certificate must certify the \emph{rank} $n$, not merely the two primes.

\subsection{Named defects in the underlying corpus}
\label{sec:defects}

The body of work this paper draws on was subjected to an adversarial review that found two
unrepaired blocking defects. Neither touches $\Fb$, and both are disclosed here rather than
smoothed.

\begin{enumerate}
\item \textbf{A vocabulary collision on the ``governing record index''.} The phrase carries
  three inequivalent meanings across the corpus, and the same statistic was reported under
  all three: (i) $m(n) = $ the most recent maximal-gap index $\le n$; (ii)
  $m(n) := \min\{m : g_m \ge g_n\}$; (iii) $T_m \le T_n$ for \emph{all} $m < n$ straddling a
  record gap. These are strictly ordered, (iii) $\Rightarrow$ (i) $\Rightarrow$ (ii), and the
  search pruning consumes only the weakest, (ii). With that named, two apparently
  contradictory findings cease to contradict: the margin of predicate~(i) decays like
  $p^{-0.83}$ ($+1.05\cdot10^{-2}$ at $3\cdot10^6$ to $+2.93\cdot10^{-6}$ at $10^{11}$, zero
  exceptions throughout), while the margin of predicate~(ii) is flat at $+0.4845$ with its
  global minimum at $n = 1879$, never approached again across eleven decades, also with zero
  exceptions. \textbf{This paper is written entirely in the disambiguated vocabulary}, but the
  upstream artifacts still carry the collision, and an inference drawn there --- that a
  \texttt{float64} check of the pruning predicate hits its noise floor at the published
  frontier --- is derived from predicate~(i) and does \emph{not} apply to predicate~(ii),
  which is $\sim10^{12}$ ulps clear of zero at every decade measured.
\item \textbf{A mis-derived bound in the Axler-sharpened version of Theorem~\ref{thm:C}.} As
  printed upstream, the theorem is false by a factor $\approx 38$ over part of its own
  validity range, because the stated upper bound on $T_m$ does not follow from its stated
  justification; the in-run numerical check reproduced the error rather than catching it. The
  review confirmed that the \emph{conclusion} is nonetheless true under the corrected tight
  form, and the repair is a one-line restatement. \textbf{It has not been applied.} This paper
  therefore states only the Dusart-only Theorem~\ref{thm:C}, whose constant $0.0623$ is
  unaffected and was independently confirmed conservative; see Remark~\ref{rem:C-b}.
\end{enumerate}

Consequently the pre-publication evidence gate on the underlying corpus stands at
\textbf{BLOCKED}, with the adversarial-review leg named as the failing one. The formal
(kernel) leg passes outright, and the adversarial-corpus leg passes ($27/27$). Nothing in
\S\S\ref{sec:setup}--\ref{sec:tension} rests on either defect.

\subsection{Source provenance, and the citation audit that has not run}
\label{sec:provenance}

Every citation in this paper traces to a row of a source ledger built from scratch for this
study: $20$ rows, of which $11$ were fetched and read at the recorded locator, $3$ were
fetched at a pagination that could not be independently confirmed, $4$ are attested at a
specific locator by two or more independent fetched primary sources, and $2$ by one such
source. No citation rests on recall. Seven source PDFs were fetched and read in full with
their MD5 digests recorded.

Two entries in \texttt{references.bib} --- \cite{carneiro2019fourier} and
\cite{visser2018andrica} --- are \textbf{pending ledger rows}: they were fetched,
version-pinned and read at the locator during this study, and proposed as ledger additions,
but were not folded into the ledger artifact. They are flagged here and in
\texttt{authoring-log.md}. Locators for \cite{carneiro2019fourier} are to
arXiv:1708.04122v2 and cannot presently be mapped to the journal's numbering.

\textbf{The citation audit for this paper has not been run, and no citation clearance is
claimed.} Three unopened or under-located sources are load-bearing, in priority order:

\begin{enumerate}
\item \cite{axler2014newbounds} --- \emph{not opened}. Quoted through Kourbatov's proofs;
  Kourbatov's Theorems 1, 3 and 5 depend on it, and so do Lemma~\ref{lem:S2},
  Theorem~\ref{thm:B} and Corollary~\ref{cor:D2}. Axler's own corrigendum moved the validity
  range of Corollary~3.5 by nine orders of magnitude. \textbf{Priority~1.}
\item \cite{granville1995cramer} --- fetched and read, but at \emph{preprint} pagination
  (pp.~1--16) against the journal's pp.~12--28. Every locator must be re-expressed against
  the journal copy. This is the load-bearing citation of the entire refutation-side argument
  (\S\ref{sec:tension}). \textbf{Priority~2.}
\item \cite{oliveira2014goldbach} --- \emph{not opened} (the publisher returned HTTP 403). It
  is the first-occurrence gap table on which Theorem~\ref{thm:verified} rests, and therefore
  the single largest external dependency of every statement about $2^{64}$ in this paper. If
  that table is incomplete, Proposition~\ref{prop:1920} and every use of
  Theorem~\ref{thm:verified} say nothing. \textbf{Priority~3.}
\end{enumerate}

Further, \cite{ribenboim2004little} (the origin locator, p.~185) was not opened and rests on
two independent citers; \cite{shanks1964maximal} and \cite{farhadian2017new} were not opened
and are cited for attribution only, their content being available at first hand elsewhere in
the bibliography; and \cite{baker2001difference} and \cite{ford2016large} were read at
abstract level only --- adequate for the qualitative use both are put to, inadequate if any
internal constant were quoted, which it is not.

\subsection{Scale, stated plainly}

The in-run exhaustive sweep reached $p < 10^{11}$, which is $8.27$ decades short of $2^{64}$.
Several of the lemma verifications were run at the smaller scale $3\cdot10^6$, which is
$12.8$ decades short of $2^{64}$. \textbf{Both sweeps verify the lemmas, not the range.}
They must never be cited as a verification of $\Fb$. The distance between what this study owns
and the published frontier is now exactly one external table and one inequality to check
against it: for every even $g$ occurring as a prime gap in $[10^{11}, 2^{64}]$, is
$q(g) > S^\sharp(g)$? The checker is written and validated on data we own; supply the table
and the frontier is re-derived in seconds.

Finally: ``unconditional'' means ``no unproved hypothesis'', not ``no computation''.
Proposition~\ref{prop:closure} proves the computation cannot be removed. Any sentence of the
form ``$\Fb$ is unconditionally true below $X$'' must, to be honest, carry the table with it.

\subsection{Dead ends, recorded so they are not rediscovered}

\begin{itemize}
\item \textbf{Littlewood oscillation is irrelevant} to the threshold, its contribution
  being $O\bigl(L^2\log\log\log x/\sqrt{x}\bigr)$, which tends to zero.
\item \textbf{Bertrand's postulate is useless} here: it implies $\Fb$ only at $n = 1$.
\item \textbf{The two-sided $\pi$-bound route} to the maximal-gap reduction cannot work
  (\S\ref{sec:ffm-counter}).
\item \textbf{More sieving is not more evidence} (\S\ref{sec:artefact}(3)). The record that
  would matter is $4.2$ decades above anything reached here.
\end{itemize}

% ==========================================================================
\section{Conclusion}
\label{sec:conclusion}

Firoozbakht's conjecture was neither proved nor refuted here. What was produced instead is a
machine-checked reduction of it to a prime-gap inequality; an unconditional, table-driven
verification architecture that reproduces the published $2^{64}$ frontier from first
principles and rests on a single explicit estimate for $\pi(x)$; an exhaustive independent
sweep to $10^{11}$ finding no counterexample and no near-miss; four theorems closing the
Riemann-Hypothesis route as a route; and the discharge, without proving it, of the obligation
this attack had ranked first --- offset by two named, unrepaired defects in the underlying
corpus and by a citation audit that has not run.

The most tractable genuinely open analytic node isolated here is the residual window of
Theorem~\ref{thm:C}: an unconditional short-interval prime count sharp to a factor
$1 + 2/\log p$, where Brun--Titchmarsh gives only a factor $2$
(Proposition~\ref{prop:obstruction}). It is stated in closed form. It is also still an open
problem in analytic number theory, not a lookup.

The standing instruction with which this work was carried out is the one it ends on. \emph{The
conjecture is open. Do not write ``Firoozbakht is true''. Do not write ``Firoozbakht is
false''. The Cram\'er--Granville tension is evidence about which way to bet, and nothing
more.}

\section*{Acknowledgements and provenance}
\addcontentsline{toc}{section}{Acknowledgements and provenance}

This paper was produced by \textbf{Noogram} as the editorial artifact of a structured,
multi-leg attack conducted on 25 July 2026. The legs contributing results reported here are:
a decomposition leg; a source-ledger leg; a concept-card leg; three proof-attempt legs; three
computational-notebook legs; Lean skeleton and probe legs; an adversarial red-team corpus leg;
an adversarial review (skeptic) leg; an evidence gate; and a synthesis leg. The delivery
posture of this artifact is \textbf{staged}: it reports a corpus whose evidence gate stands at
BLOCKED (\S\ref{sec:defects}) and whose citation audit has not run
(\S\ref{sec:provenance}). It is not a seal.

\printbibliography[heading=bibintoc,title={References}]

\end{document}
